\documentclass[12pt,a4paper]{article}
\usepackage{fontspec}
\usepackage{amsmath, amssymb, amstext}
\usepackage{graphicx}
\usepackage{geometry}
\usepackage{xcolor}
\usepackage{fancyhdr}
\geometry{margin=1in}
\pagestyle{fancy}
\setlength{\headheight}{15pt} 
\lhead{کوییز درس فیزیک مدرن}
\chead{دانشکده فیزیک دانشگاه صنعتی خواجه نصیرالدین طوسی}
\rhead{\today}
\usepackage{xepersian}
\settextfont{B Nazanin}

\begin{document}
	
	\begin{tabular}{|p{7cm}|p{7cm}|}
		\hline
		\textbf{نام و نام خانوادگی:} & \textbf{شماره دانشجویی:} \\
		\hline
	\end{tabular}
	
	
	\section*{سوالات:}
	\begin{enumerate}
		\item 
		
		اتمی با جرم $m_1 = m$ در جهت مثبت $x$ با سرعت $v_1 = v$ حرکت می‌کند. این اتم با اتمی به جرم $m_2 = 2m$ که در جهت مثبت $y$ با سرعت $v_2 = 2v/3$ حرکت می‌کند، برخورد کرده و به آن می‌چسبد. سرعت حاصل و جهت حرکت ترکیب را پیدا کنید، و همچنین انرژی جنبشی تلف شده در این برخورد ناکشسان را بیابید.
		\item 
		
		تابع چگالی احتمال سرعت برای یک مولکول گاز ایده‌ال با جرم $m$ در دمای $T$ به صورت زیر است:
		
		\[
		f(v) = \sqrt{\frac{2}{\pi}} \, \left( \frac{m}{k_B T} \right)^{3/2} v^{2} \, \exp\left( -\frac{m v^{2}}{2 k_B T} \right), \quad v \ge 0
		\]
		
		که $k_B$ ثابت بولتزمن است.
		\\
		\\
		
		ثابت کنید:
		
		\subsection*{الف) نرمالیزه بودن}
		\[
		\int_{0}^{\infty} f(v) \, dv = 1
		\]
		
		\subsection*{ب) ممان اول (میانگین سرعت)}
		\[
		\langle v \rangle = \int_{0}^{\infty} v \, f(v) \, dv = \sqrt{\frac{8 k_B T}{\pi m}}
		\]
		
		\subsection*{پ) ممان دوم (میانگین مربع سرعت)}
		\[
		\langle v^{2} \rangle = \int_{0}^{\infty} v^{2} \, f(v) \, dv = \frac{3 k_B T}{m}
		\]
		
		\subsection*{ت) واریانس سرعت}
		\[
		\sigma_v^{2} = \langle v^{2} \rangle - \langle v \rangle^{2}
		= \frac{3 k_B T}{m} - \frac{8 k_B T}{\pi m}
		= \frac{k_B T}{m} \left( 3 - \frac{8}{\pi} \right)
		\]
		
	
		
	\end{enumerate}
	
	\vspace{10pt}
	\textbf{موفق باشید.}
	
\end{document}
